\documentclass[a4paper]{memoir}
\usepackage{graphicx}
%% Language and font encodings
\usepackage{microtype}
\usepackage[paperheight=88.9mm,paperwidth=63.5mm,margin=5mm]{geometry}
\usepackage{ragged2e}
\usepackage{adjustbox}
\usepackage{pgffor}
\usepackage{emoji}
\usepackage{tikz}
\newcommand{\cardtext}{1-$\infty$}
\newcommand{\grouptext}{2-6}
\newcommand{\agetext}{6+}
\newcommand{\clocktext}{20$''$}
\newcommand{\languagecrop}{english_crop}
\newcommand{\colorcount}{6}
\newcommand{\cardtitle}{}
\def\colors{0/myyellow, 1/myorange, 2/myred, 3/mypurple, 4/myblue, 5/mygreen}
\begin{document}
{\footnotesize
\input{../../../rule_lib/rule_lib.tex}
\noindent
\textbf{\emoji{party-popper} Unomito Variant.} Play with the classic Unomito rules \textbf{or} the 1pile variant, now including additional cards: \textbf{the \infty~cards}.
\\
\\
\noindent
\textbf{\emoji{game-die} The \infty~is a wild card.} These 3 extra cards act as wild cards, letting you play them without the risk of being accused of lying.
\newpage
\noindent
\textbf{\emoji{party-popper} Classic Unomito.} When an \infty~is revealed face up, the player who played it can choose its color and number.
\\
\\
\noindent
\textbf{\emoji{party-popper} Unomito 1pile.} When an \infty~is revealed face up, it takes the color and number declared by the player when it was played face down.
}\end{document}
